\section{Introduction}

Mathematical modeling of infectious diseases has become an indispensable tool for public health decision-making and outbreak forecasting.
Traditional compartmental models, such as the Susceptible-Infected-Recovered (SIR) framework, have historically relied on the assumption of homogeneous mixing, where every individual in a population has an equal probability of contacting any other individual.
However, recent empirical evidence and the global experience of the COVID-19 pandemic have clearly demonstrated that human interaction patterns are far from uniform.
\textcite{britton2025improving} argue that neglecting the inherent heterogeneity in social contacts can lead to significant biases in estimating the basic reproduction number ($R_0$) and the final epidemic size.
Therefore, incorporating realistic contact structures through network-based modeling is essential for accurately capturing the dynamics of disease transmission.

The foundation for constructing such realistic models lies in the empirical study of social mixing patterns.
The landmark POLYMOD study provided the first large-scale, population-based dataset on social contacts across eight European countries, revealing consistent patterns of age-dependent mixing and household clustering \parencite{mossong2008social}.
Since then, these contact matrices have been widely used to project transmission dynamics in various global settings, including low- and middle-income countries \parencite{prem2017projecting}.
Despite the richness of these survey-based data, a significant challenge remains in "reverse-engineering" these observed contact frequencies into a generative network model that can be used for high resolution stochastic simulations.

From a statistical perspective, inferring the underlying parameters of a contact network from survey data presents a formidable challenge.
Because the likelihood function for complex network-based epidemic models is often analytically intractable or computationally prohibitive, standard Maximum Likelihood Estimation (MLE) is frequently inapplicable \parencite{mckinley2009inference}.
Recent advancements in computational statistics, particularly Likelihood-Free Inference (LFI) and Synthetic Likelihood methods, have opened new avenues for parameter estimation in high-dimensional settings \parencite{ong2018likelihoodfree}.
By leveraging summary statistics that represent key network properties, these methods allow for the robust calibration of generative models without the need for an explicit likelihood function \parencite{picchini2023sequentially}.

The primary objective of this thesis is to develop a robust inference framework that reproduces the contact patterns observed in the POLYMOD dataset.
Specifically, we aim to estimate the "rules" or parameters that govern human connections and validate the resulting model through a synthetic population approach.
By reconstructing a structured contact network that explicitly accounts for sociability variance and school-based cohorts, we evaluate the effectiveness of targeted public health interventions, such as school closures.
This research bridges the gap between empirical contact surveys and generative network modeling, providing a scalable foundation for simulating realistic outbreak scenarios and evaluating the non-linear impacts of behavioral shifts.

The remainder of this thesis is organized as follows.
Section 2 describes the POLYMOD dataset and the data cleaning process.
Section 3 details the formulation of the network model and the Synthetic Likelihood MCMC inference pipeline.
Section 4 presents the results of the epidemic simulations and compares the effectiveness of different intervention strategies.
Finally, Section 5 discusses the implications of our findings, acknowledges limitations such as contact substitution, and outlines directions for future work.